% Einheit 4 – Graph und Ableitungsgraph
\setcounter{aufgabe}{28}
\einheitenkopf[e4][4 Graph und Ableitungsgraph]{Einheit 4 von 5 · Graph und Ableitungsgraph}
\zweigzeile{Hier lernst du, vom Graphen von $f$ auf den Graphen von $f'$ zu schließen und umgekehrt und Graphen aus Eigenschaften zu skizzieren · kennst du seit Q1 oder Q2 – je nach Buch · Abitur GK · baut auf: Einheit 1 bis 3, Ableitung als Steigung (Nr. 8)}

\begin{aufgabe}{Ich kann am Graphen von $f$ ablesen, wo $f'$ positiv, negativ oder null ist.}
Lies am Graphen von $f$ ab (im gezeichneten Bereich).

\begin{ksys}[xmin=-3,xmax=3,ymin=-4,ymax=4,ablesen]
  \funktion{0.25*\x^4-2*\x^2+1}{f}
\end{ksys}

\begin{teile}
\teil Stellen, an denen $f'(x) = 0$ ist: \leerfeld
\teil Bereiche, in denen $f'(x) > 0$ ist: \leerfeld
\teil Bereiche, in denen $f'(x) < 0$ ist: \leerfeld
\teil Ist $f'(1)$ positiv oder negativ? \leerfeld
\anweisung{Markiere im Graphen und begründe.}
\teil Markiere den Wendepunkt, in dem der Graph fällt. Begründe, dass $f'$ dort negativ ist und dort seinen kleinsten Wert im Bereich $0 \le x \le 2$ hat.
\end{teile}
\end{aufgabe}

\begin{aufgabe}{Ich kann am Graphen von $f'$ ablesen, was der Graph von $f$ macht.}
Die Abbildung zeigt den Graphen der Ableitung $f'$ einer Funktion $f$.

\begin{ksys}[xmin=-2,xmax=4,ymin=-3,ymax=3,ablesen]
  \funktion{0.5*\x^2-\x-1.5}{f'}
\end{ksys}

\begin{teile}
\teil Lies $f'(1)$ ab. \leerfeld
\teil An welcher Stelle hat der Graph von $f$ eine Wendestelle? \leerfeld
\anweisung{Gib an und begründe mit dem Graphen von $f'$.}
\teil An welchen Stellen hat $f$ Extremstellen, und welcher Art sind sie?
\teil In welchen Bereichen steigt der Graph von $f$?
\end{teile}
\end{aufgabe}

\begin{aufgabe}{Ich kann die Graphen von $f$ und $f'$ aus berechneten Werten zeichnen.}
Gegeben ist $f(x) = 0{,}5x^3 - 1{,}5x^2$.
\begin{teile}
\teil Berechne $f(x)$ und $f'(x)$ und trage sie in die Tabellen ein.
\end{teile}

\wertetabelle{x}{f(x)}{-1,0,1,2,3}

\wertetabelle{x}{f'(x)}{-1,0,1,2,3}

\begin{teile}
\teil Zeichne den Graphen von $f$ in das obere und den Graphen von $f'$ in das untere Koordinatensystem.
\teil Markiere die Punkte des Graphen von $f$, die zu den Nullstellen von $f'$ gehören. Was für Punkte sind das?
\end{teile}

\ableitungspaar[xmin=-2,xmax=4,ymin=-3,ymax=2,ymin2=-2,ymax2=5,karo=0.8]{}{f}{}{f'}

\end{aufgabe}

\begin{aufgabe}{Ich kann einen möglichen Graphen zu gegebenen Eigenschaften skizzieren.}
Skizziere einen möglichen Graphen von $f$ in das Koordinatensystem mit demselben Buchstaben.
\begin{teile}
\teil $f$ hat den Hochpunkt $H(-1\,|\,3)$ und den Tiefpunkt $T(2\,|\,{-}2)$; für $x \to \infty$ gilt $f(x) \to \infty$.
\teil $f$ hat den Sattelpunkt $S(0\,|\,2)$ und den Tiefpunkt $T(2\,|\,{-}1)$; für $x \to \pm\infty$ gilt $f(x) \to \infty$.
\anweisung{Begründe.}
\teil $f$ in b) ist ganzrational. Welchen Grad hat $f$ mindestens?
\end{teile}

\begin{ksys}[xmin=-3,xmax=4,ymin=-3,ymax=4]
  \node[anchor=north west,font=\bfseries] at (-3,4) {a)};
\end{ksys}\ksysabstand\begin{ksys}[xmin=-3,xmax=4,ymin=-3,ymax=4]
  \node[anchor=north west,font=\bfseries] at (-3,4) {b)};
\end{ksys}

\end{aufgabe}

\begin{aufgabe}{Ich kann Aussagen über Graphen und Ableitungen beurteilen.}
Die Abbildung zeigt den Graphen der Differenzfunktion $d$ mit $d(x) = f(x) - g(x)$ zweier Funktionen $f$ und $g$. Beurteile die Aussage.

\begin{ksys}[xmin=0,xmax=5,ymin=-4,ymax=3,ablesen]
  \funktion{-\x^2+5*\x-4}{d}
\end{ksys}

\begin{teile}
\teil „Für $1 < x < 4$ liegt der Graph von $f$ oberhalb des Graphen von $g$.“
\teil „An der Stelle $x = 2{,}5$ haben die Graphen von $f$ und $g$ dieselbe Steigung.“
\anweisung{Beurteile die Aussagen rechnerisch.}
\teil Gegeben ist $f(x) = x^3 - 6x^2 + 10x$. (1) „Die Normale im Wendepunkt schneidet die $y$-Achse bei $y = 3$.“ (2) „Es gibt eine Stelle, an der der Graph von $f$ die Steigung $-3$ hat.“ (Abitur 2024 GK)
\end{teile}
\end{aufgabe}

\begin{aufgabe}{Ich finde den Fehler beim Lesen von Graph und Ableitungsgraph.}
Finde den Fehler, benenne ihn und korrigiere die Aussage.
\begin{teile}
\teil Der Graph von $f'$ hat bei $x = 1$ einen Tiefpunkt. Lisa schreibt: „Also hat $f$ bei $x = 1$ einen Tiefpunkt.“
\teil Max liest am Graphen von $f$ den Punkt $P(2\,|\,3)$ ab und schreibt: „$f'(2) = 3$.“
\end{teile}
\end{aufgabe}

\begin{aufgabe}{Ich kann Zusammenhänge zwischen $f$ und $f'$ begründen.}
Begründe ohne Rechnung.
\begin{teile}
\teil Warum liegen die Nullstellen von $f'$ genau unter den Hoch- und Tiefpunkten von $f$?
\teil Warum ist der Graph von $f'$ bei einer ganzrationalen Funktion $f$ vom Grad drei eine Parabel?
\end{teile}
\end{aufgabe}

\begin{aufgabe}{Ich kann aus dem Graphen einer Änderungsrate auf den Bestand schließen.}
Die Abbildung zeigt die Zuflussrate $r$ eines Wasserbeckens in m$^3$ pro Stunde ($t$ in Stunden nach Beobachtungsbeginn, $0 \le t \le 8$); negative Werte bedeuten, dass Wasser abfließt.

\begin{ksys}[xmin=0,xmax=8,ymin=-4,ymax=5,ablesen,xlabel=$t$,ylabel=$r(t)$]
  \funktion{-0.5*\x^2+4*\x-3.5}{r}
\end{ksys}

\begin{teile}
\teil In welchem Zeitraum nimmt die Wassermenge im Becken zu?
\teil Zu welchem Zeitpunkt ist am meisten Wasser im Becken? Begründe.
\teil Zu welchem Zeitpunkt nimmt die Wassermenge am schnellsten zu?
\end{teile}
\end{aufgabe}
